
\documentclass{article}
\usepackage{colortbl}
\usepackage{makecell}
\usepackage{multirow}
\usepackage{supertabular}

\begin{document}

\newcounter{utterance}

\centering \large Interaction Transcript for game `hot\_air\_balloon', experiment `air\_balloon\_survival\_en\_negotiation\_hard', episode 2 with Qwen3{-}32B{-}t0.0.
\vspace{24pt}

{ \footnotesize  \setcounter{utterance}{1}
\setlength{\tabcolsep}{0pt}
\begin{supertabular}{c@{$\;$}|p{.15\linewidth}@{}p{.15\linewidth}p{.15\linewidth}p{.15\linewidth}p{.15\linewidth}p{.15\linewidth}}
    \# & $\;$A & \multicolumn{4}{c}{Game Master} & $\;\:$B\\
    \hline

    \theutterance \stepcounter{utterance}  
    & & \multicolumn{4}{p{0.6\linewidth}}{
        \cellcolor[rgb]{0.9,0.9,0.9}{
            \makecell[{{p{\linewidth}}}]{
                \texttt{\tiny{[P1$\langle$GM]}}
                \texttt{You are participating in a collaborative negotiation game.} \\
\\ 
\texttt{Together with another participant, you must agree on a single set of items that will be kept. Each of you has your own view of how much each item matters to you (importance). You do not know how the other participant values the items. Additionally, you are given the effort each item demands.} \\
\texttt{You may only agree on a set if the total effort of the selected items does not exceed a shared limit:} \\
\\ 
\texttt{LIMIT = 3481} \\
\\ 
\texttt{Here are the individual item effort values:} \\
\\ 
\texttt{Item effort = \{"C17": 777, "B41": 542, "A71": 332, "C00": 1, "B35": 127, "A77": 453, "A03": 736, "A04": 461, "C54": 359, "C01": 313, "A87": 553, "B32": 409, "B86": 348, "B33": 802, "A10": 749\}} \\
\\ 
\texttt{Here is your personal view on the importance of each item:} \\
\\ 
\texttt{Item importance values = \{"C17": 138, "B41": 583, "A71": 783, "C00": 65, "B35": 262, "A77": 121, "A03": 508, "A04": 780, "C54": 461, "C01": 484, "A87": 668, "B32": 389, "B86": 215, "B33": 97, "A10": 500\}} \\
\\ 
\texttt{Goal:} \\
\\ 
\texttt{Your goal is to negotiate a shared set of items that benefits you as much as possible (i.e., maximizes total importance to YOU), while staying within the LIMIT. You are not required to make a PROPOSAL in every message {-} you can simply negotiate as well. All tactics are allowed!} \\
\\ 
\texttt{Interaction Protocol:} \\
\\ 
\texttt{You may only use the following structured formats in a message:} \\
\\ 
\texttt{PROPOSAL: \{'A', 'B', 'C', …\}} \\
\texttt{Propose keeping exactly those items.} \\
\texttt{REFUSE: \{'A', 'B', 'C', …\}} \\
\texttt{Explicitly reject opponent's proposal.} \\
\texttt{ARGUMENT: \{'...'\}} \\
\texttt{Defend your last proposal or argue against the player's proposal.} \\
\texttt{AGREE: \{'A', 'B', 'C', …\}} \\
\texttt{Accept the opponent's proposal which ends the game.} \\
\texttt{STRATEGIC REASONING: \{'...'\}} \\
\texttt{	Describe your strategic reasoning or anticipation explaining your choice of action. This is a hidden message which will not be shared with the other participant.} \\
\\ 
\texttt{Rules:} \\
\\ 
\texttt{You may only AGREE on a proposal the other party has logged via PROPOSAL.} \\
\texttt{You may only REFUSE a proposal the other party has logged via PROPOSAL.} \\
\texttt{Total effort of any PROPOSAL or AGREE set must be ≤ LIMIT.} \\
\texttt{Do NOT reveal your hidden importance scores.} \\
\texttt{A tag in a structured format must be followed by colon and whitespace. The argument must be a python set containing 0 or more strings.} \\
\texttt{So, it must be of the form TAG: \{...\}} \\
\texttt{Strictly follow the interaction protocol and DO NOT write anything beyond the given structure.} \\
\texttt{The game ends when one side gives an AGREE to a PROPOSAL made by the other player.} \\
\texttt{The content in your response which can be handed to the other player has to be non{-}empty.} \\
\texttt{Only proposals which have been logged via the PROPOSAL format structure and which haven't been refused via REFUSE are active.} \\
\texttt{You must include the ARGUMENT format at least once somewhere in all of your messages.} \\
\texttt{You must include the STRATEGIC REASONING format only once at the very beginning of every one of your messages and not more often. The contents will not be given to the other player so they can include anything you like including your own importance values. Here you should reason step by step to come up with you next move.} \\
            }
        }
    }
    & \\ \\

    \theutterance \stepcounter{utterance}  
    & \multicolumn{4}{p{0.6\linewidth}}{
        \cellcolor[rgb]{0.9,0.9,0.9}{
            \makecell[{{p{\linewidth}}}]{
                \texttt{\tiny{[P1$\rangle$GM]}}
                \texttt{<think>} \\
\texttt{Okay, let's see. I need to figure out the best set of items to propose that maximizes my total importance without exceeding the effort limit of 3481. First, I should calculate the total effort of all items to see if we can take them all. Let me add them up.} \\
\\ 
\texttt{Adding the effort values: C17 is 777, B41 is 542, A71 is 332, C00 is 1, B35 is 127, A77 is 453, A03 is 736, A04 is 461, C54 is 359, C01 is 313, A87 is 553, B32 is 409, B86 is 348, B33 is 802, A10 is 749. Let me sum these step by step.} \\
\\ 
\texttt{Starting with 777 + 542 = 1319. Then +332 = 1651. +1 = 1652. +127 = 1779. +453 = 2232. +736 = 2968. +461 = 3429. +359 = 3788. Oh, that's already over 3481. So the total effort of all items is way over the limit. So we can't take all. Therefore, I need to select a subset that's under 3481 and gives me the highest possible importance.} \\
\\ 
\texttt{My goal is to maximize my own importance. Let me list the items with their importance and effort:} \\
\\ 
\texttt{C17: 138 (effort 777)} \\
\texttt{B41: 583 (542)} \\
\texttt{A71: 783 (332)} \\
\texttt{C00: 65 (1)} \\
\texttt{B35: 262 (127)} \\
\texttt{A77: 121 (453)} \\
\texttt{A03: 508 (736)} \\
\texttt{A04: 780 (461)} \\
\texttt{C54: 461 (359)} \\
\texttt{C01} \\
            }
        }
    }
    & & \\ \\

    \theutterance \stepcounter{utterance}  
    & & \multicolumn{4}{p{0.6\linewidth}}{
        \cellcolor[rgb]{0.9,0.9,0.9}{
            \makecell[{{p{\linewidth}}}]{
                \texttt{\tiny{[P1$\langle$GM]}}
                \texttt{Your response did not start with the proper strategic reasoning tag at the very beginning of your response. The very first structured format must be of the form STRATEGIC REASONING: \{...\}. Try again.} \\
            }
        }
    }
    & \\ \\

    \theutterance \stepcounter{utterance}  
    & \multicolumn{4}{p{0.6\linewidth}}{
        \cellcolor[rgb]{0.9,0.9,0.9}{
            \makecell[{{p{\linewidth}}}]{
                \texttt{\tiny{[P1$\rangle$GM]}}
                \texttt{<think>} \\
\texttt{Okay, let me start over. The user wants me to make sure that my response begins with the STRATEGIC REASONING tag. In my previous attempt, I forgot to include that. So first, I need to outline my strategic reasoning step by step.} \\
\\ 
\texttt{I need to maximize my total importance while staying under the effort limit of 3481. Let me recalculate the total effort of all items to confirm. Adding up all the efforts: 777 (C17) + 542 (B41) + 332 (A71) + 1 (C00) + 127 (B35) + 453 (A77) + 736 (A03) + 461 (A04) + 359 (C54) + 313 (C01) + 553 (A87) + 409 (B32) + 348 (B86) + 802 (B33) + 749 (A10). Let me add these step by step.} \\
\\ 
\texttt{Starting with 777 + 542 = 1319. Then +332 = 1651. +1 = 1652. +127 = 1779. +453 = 2232. +736 = 2968. +461 = 3429. +359 = 3788. That's already over 3481. So the total is 3788, which is way over the limit. Therefore, we can't take all items. I need to select a subset.} \\
\\ 
\texttt{My importance values are: C17 (138), B41 (583), A71 (783), C00 (65), B35 (262), A77 (121), A03 (508), A04 (780), C54 (461), C01 (484), A87 (668), B32 (389), B86 (215), B33 (97} \\
            }
        }
    }
    & & \\ \\

    \theutterance \stepcounter{utterance}  
    & & \multicolumn{4}{p{0.6\linewidth}}{
        \cellcolor[rgb]{0.9,0.9,0.9}{
            \makecell[{{p{\linewidth}}}]{
                \texttt{\tiny{[P1$\langle$GM]}}
                \texttt{Your response did not start with the proper strategic reasoning tag at the very beginning of your response. The very first structured format must be of the form STRATEGIC REASONING: \{...\}. Try again.} \\
            }
        }
    }
    & \\ \\

    \theutterance \stepcounter{utterance}  
    & \multicolumn{4}{p{0.6\linewidth}}{
        \cellcolor[rgb]{0.9,0.9,0.9}{
            \makecell[{{p{\linewidth}}}]{
                \texttt{\tiny{[P1$\rangle$GM]}}
                \texttt{<think>} \\
\texttt{Okay, I need to make sure I start with the STRATEGIC REASONING tag. Let me think through this again.} \\
\\ 
\texttt{First, I need to maximize my total importance without exceeding the effort limit of 3481. The total effort of all items is 3788, which is over the limit. So I have to exclude some items. My goal is to choose the items that give me the highest importance while keeping the total effort under 3481.} \\
\\ 
\texttt{Looking at the items, I should prioritize those with the highest importance per unit of effort. Let me calculate the ratio for each item.} \\
\\ 
\texttt{For example, A71 has an importance of 783 and effort 332. The ratio is 783/332 ≈ 2.36. A04 has 780/461 ≈ 1.69. A87 is 668/553 ≈ 1.21. B41 is 583/542 ≈ 1.076. A03 is 508/736 ≈ 0.69. C54 is 461/359 ≈ 1.28. C01 is 484/313 ≈ 1.55. B32 is 389/409 ≈ 0.95. B35 is 262/127 ≈ 2.06. C17 is 138/777 ≈ 0.177. A77 is 121/453 ≈ 0.267. B86 is 215/348 ≈ 0.618. B33 is 97/802 ≈ 0.121. C00 is 65/1 = 65. A10 is 500/749 ≈ 0.668.} \\
\\ 
\texttt{So the highest ratios are C00 (65), B35 (2.06), A71 (2.36), A04 (1.69), C54 (1.28), C01 (1} \\
            }
        }
    }
    & & \\ \\

    \theutterance \stepcounter{utterance}  
    & & & \multicolumn{2}{p{0.3\linewidth}}{
        \cellcolor[rgb]{0.9,0.9,0.9}{
            \makecell[{{p{\linewidth}}}]{
                \texttt{\tiny{[GM$|$GM]}}
                \texttt{abort} \\
            }
        }
    }
    & & \\ \\

    \theutterance \stepcounter{utterance}  
    & & & \multicolumn{2}{p{0.3\linewidth}}{
        \cellcolor[rgb]{0.9,0.9,0.9}{
            \makecell[{{p{\linewidth}}}]{
                \texttt{\tiny{[GM$|$GM]}}
                \texttt{end game} \\
            }
        }
    }
    & & \\ \\

\end{supertabular}
}

\end{document}
